\section{Processament asíncron i cues de treball}

Una cua de treball separa la producció d'una tasca de la seva execució. El
productor crea una unitat de treball amb les dades necessàries i la cua la manté
fins que un worker disposa de capacitat per processar-la. Aquest esquema permet
regular pics de demanda, distribuir feina entre diversos processos i tractar
fallades sense bloquejar necessàriament la petició que ha originat la tasca. El
model bàsic es pot descriure amb l'instant d'arribada, la durada del servei, la
prioritat, l'estat i el worker assignat a cada tasca.

FIFO (\emph{first in, first out}) és la disciplina de partida més directa: les
tasques es retiren segons l'ordre d'arribada. Tanmateix, l'ordre que observa
l'aplicació no sempre coincideix amb una seqüència FIFO perfecta. En una cua
operativa hi poden intervenir diversos consumidors, tasques lliurades però no
confirmades, reintents o prioritats. RabbitMQ documenta que aquests factors
poden modificar l'ordre efectiu de lliurament \parencite{rabbitmqPriority}. Una
simulació ha de declarar quins efectes representa i quins simplifica.

BullMQ mostra una implementació del mateix problema en l'ecosistema Node.js.
Les tasques prioritzades s'ordenen pel valor assignat i conserven FIFO dins del
mateix nivell; la inserció prioritzada té un cost logarítmic respecte del nombre
de tasques del conjunt ordenat \parencite{bullmqPrioritized}. Aquest detall és
rellevant per entendre una eina real, però no determina el resultat del
simulador: el TFM compararà les conseqüències de cada política sobre la càrrega
modelada, no el rendiment intern d'una biblioteca concreta.

La concurrència obliga a precisar el moment de la decisió. En un model no
preemptiu, una tasca que ja s'està executant manté el worker fins que finalitza;
la política només escull entre les tasques pendents quan queda capacitat lliure.
Amb diversos workers, dues tasques consecutives poden acabar en un ordre
diferent si tenen durades diferents. L'ordre de selecció, l'ordre d'inici i
l'ordre de finalització són, per tant, conceptes separats.

La plataforma treballarà amb una abstracció més simple que un broker productiu.
No reproduirà protocols de confirmació, replicació ni persistència distribuïda.
Representarà els elements que afecten directament la pregunta experimental:
arribades, durades, prioritats, workers i, si entra dins de l'abast final,
fallades i reintents. Aquesta delimitació permet atribuir els canvis observats
a la política en lloc de barrejar-los amb decisions d'infraestructura.
